\subsection{Formatting HTML inside Java}\label{sec:FormattingHTMLInsideJava}
Embedded HTML\index{HTML} snippets are usually used to format the output of a program or an application, not only in the context of web applications.

Below some recommendations that were translated from \autocite{selfHTML:GuterHTMLStil}:
\begin{itemize}
	\item{Your HTML markup should be error-free (valid) so that it can be read by browsers and other parsers, such as screen readers. HTML\index{HTML} can only be validated if you include a \verb#doctype# declaration\footnote{This is only relevant if you embedd a complete HTML document into your Java code, and not just HTML snippets.}.} 
	\item{Virtues borrowed from XML\index{XML}: 
		\begin{itemize}
			\item{Well-formed syntax improves readability and enables processing as XHTML\index{XHTML}.}
			\item{All element names, attribute names, and their values are written in lowercase. Attribute values are always enclosed in doubel quotation marks.}
			\item{Use \bdq{meaningful} class names that describe function rather than visual styling.} 
			\item{Elements that do not strictly require a closing tag in HTML (e.g., \verb#p#, \verb#th#, \verb#td#, \verb#dt#, \verb#dd#, \verb#li#) should always be written with a closing tag.}
			\item{Use indentation and blank lines to structure your code clearly.}
		\end{itemize}}
\end{itemize}
\textbf{Note:} Since XHTML\index{XHTML} is an XML\index{XML} dialect, the points listed above are mandatory in XHTML\index{XHTML}, whereas in HTML\index{HTML} they are \bdq{merely} a matter of good practice.

\keeplisting{19}
This short HTML\index{HTML} document
\begin{lstlisting}[language=HTML]
<!DOCTYPE html>
<html>
    <head>
        <meta charset="UTF-8" />
        <title>Title</title>
    </head>
    <body>
        <h1>Headline</h2>
        <p>Text …</p>
        <p>More text …</p>
        <ol>
            <li>Entry 1</li>
            <li>Entry 2</li>
            <li>Entry 3</li>
        </ol>
        <p>Final text!</p>
    </body>
</html>
\end{lstlisting}

\keeplisting{21}
would be embedded like this:
\begin{lstlisting}
final var htmlDocumentString = """
````<!DOCTYPE html>
````<html>
````    <head>
````        <meta charset="UTF-8" />
````        <title>Title</title>
````    </head>
````    <body>
````        <h1>Headline</h2>
````        <p>Text …</p>
````        <p>More text …</p>
````        <ol>
````            <li>Entry 1</li>
````            <li>Entry 2</li>
````            <li>Entry 3</li>
````        </ol>
````        <p>Final text!</p>
````    </body>
````</html>
````""";
\end{lstlisting}

\keeplisting{11}
If you have to embed an HTML fragment – most probably, this would be a \verb#<p># element – this would look like this:
\begin{lstlisting}
final var htmlFragment = """
````<p>Text …</p>
````<p>More text …</p>
````    <ol>
````        <li>Entry 1</li>
````        <li>Entry 2</li>
````        <li>Entry 3</li>
````    </ol>
````<p>Final text!</p>
````""";
\end{lstlisting}

\keeplisting{4}
Or for an \verb#<li># element:
\begin{lstlisting}
final var htmlFragment = """
````<li>Entry 1</li>
````""";
\end{lstlisting}

This means that the top level element will not be indented. If you need to compose an HTML\index{HTML} document from the snippets, jyou can correct the indentation through calling \lstinline|String::indent|\index{java.lang!String!indent()} \autocite{ORACLE_DOC_STRING:indent}. This is explained in more detail in chapter \tqfullvref{sec:TextBlocks}.

%==============================================================================
% End of file!!
%==============================================================================

